Последовательность логарифмов факториалов обладает более сильным
свойством, чем плотность множества её дробных частей: она равномерно
распределена по модулю $1$. Для десятичного основания это классический
результат о ведущих цифрах факториалов~\cite{Diaconis1977};
ниже приведено доказательство для любого фиксированного основания.

\begin{lemma}\label{lem:logfact}
Пусть $b>1$. Последовательность $(\log_b(n!))_{n\ge1}$
равномерно распределена по модулю $1$. В частности, для каждого
$[\alpha,\beta)\subseteq[0,1)$ при $\alpha<\beta$ существует
бесконечно много $n$ таких, что
$\{\log_b(n!)\}\in[\alpha,\beta)$.
\end{lemma}
\begin{proof}
Зафиксируем целое $h\ne0$ и положим
\[
 u_n=\exp(2\pi\mathrm i h\log_b(n!)),\qquad
 \tau=\frac{2\pi h}{\log b}\ne0.
\]
Для каждого фиксированного целого $r\ge1$ имеем
\[
 u_{n+r}\overline{u_n}
 =\exp\left(\mathrm i\tau\sum_{j=1}^r\log(n+j)\right)
 =n^{\mathrm i\tau r}+O_{b,h,r}(n^{-1}).
\]
Здесь использовано равенство
$\log(n+j)=\log n+O_r(n^{-1})$.
Для фиксированного вещественного $a$ сравнение суммы с интегралом даёт
\[
 \sum_{n=1}^M n^{\mathrm i a}
 =\frac{M^{1+\mathrm i a}}{1+\mathrm i a}+O_a(1+\log M).
\]
Действительно, ошибка сравнения ограничена величиной
$1+\int_1^M |(x^{\mathrm i a})'|\,dx
\le1+|a|\log M$, а интеграл вычисляется непосредственно.
Следовательно, для
$C_N(r)=\sum_{n=1}^{N-r}u_{n+r}\overline{u_n}$ получаем
\[
 \limsup_{N\to\infty}\frac{|C_N(r)|}{N}
 =\frac1{\sqrt{1+\tau^2r^2}}.
\]

Приведём необходимую оценку средних, чтобы явно указать все шаги.
Продолжим $u_n$ нулями за пределы $1\le n\le N$.
Для фиксированного целого $H\ge2$ имеем
\[
 H\sum_{n=1}^N u_n
 =\sum_{m=1}^{N+H-1}\sum_{j=0}^{H-1}u_{m-j}.
\]
Неравенство Коши--Буняковского и раскрытие квадрата дают
\[
 H^2\left|\sum_{n=1}^N u_n\right|^2
 \le (N+H-1)\left(
 HN+2\sum_{r=1}^{H-1}(H-r)\operatorname{Re}C_N(r)\right).
\]
Разделив на $H^2N^2$ и сначала переходя к верхнему пределу по $N$,
получаем
\[
\begin{aligned}
 \limsup_{N\to\infty}
 \left|\frac1N\sum_{n=1}^N u_n\right|^2
 &\le \frac1H+
 \frac2H\sum_{r=1}^{H-1}
 \left(1-\frac rH\right)\frac1{\sqrt{1+\tau^2r^2}}\\
 &\le \frac1H+\frac{2}{|\tau|H}\sum_{r=1}^{H-1}\frac1r
 \xrightarrow[H\to\infty]{}0.
\end{aligned}
\]
По критерию Вейля последовательность равномерно распределена.
Доля попаданий в интервал положительной длины стремится к
$\beta-\alpha>0$, поэтому число таких попаданий бесконечно.
\end{proof}
